\begin{ouverture}[mole]{6}{Compter les entités}{Constitution et transformations}{Observer · Modéliser · Expérimenter · Expliquer}
\objectif{Calculer la masse d’une entité à partir de sa formule.}
\objectif{Passer d’une masse à un nombre d’entités.}
\objectif{Utiliser la définition de la mole.}
\prerequis{Proportionnalité, conversions et lecture de graphiques.}
\end{ouverture}
\essentiel
\essSection{1}{De la formule à la masse}
Une formule brute indique les atomes présents dans une entité et leur nombre. Dans \ce{CO2}, il y a un atome de carbone et deux d'oxygène. Pour calculer la masse d'une molécule à la précision du cours, additionner les masses des atomes qui la constituent.
\begin{exemple}[Une molécule d'eau]Avec $m_H=\qty{1,67e-27}{kg}$ et $m_O=\qty{2,66e-26}{kg}$, $m_{\ce{H2O}}=2m_H+m_O=\qty{2,994e-26}{kg}$. Les indices multiplient les masses correspondantes.\end{exemple}
\essSection{2}{Une masse contient combien d’entités ?}
\begin{loi}[Entités identiques]Pour un échantillon pur de masse $m$, le nombre d'entités vaut $N=m/m_{\mathrm{entité}}$. Les deux masses doivent être exprimées dans la même unité. $N$ est un nombre sans unité.\end{loi}
Les valeurs obtenues sont souvent immenses ; l'écriture scientifique permet de les manipuler. Dans un mélange, il faut d'abord isoler la masse de l'espèce étudiée : on ne divise pas la masse totale du mélange par la masse d'une seule sorte d'entité.
\schema[corrige]{mole}
\essSection{3}{La mole, un paquet de comptage}
\begin{definition}[Quantité de matière]Une mole contient exactement $\num{6,02214076e23}$ entités. La constante d'Avogadro vaut $N_A=\qty{6,02214076e23}{mol^{-1}}$ ; $n=N/N_A$ s'exprime en moles. En entraînement on utilise souvent $N_A\simeq\qty{6,02e23}{mol^{-1}}$.\end{definition}
Une mole d'eau et une mole de dioxyde de carbone contiennent le même nombre de molécules, mais leurs masses diffèrent. Le parcours demandé ici est : masse d'une entité, nombre d'entités, quantité de matière. La masse molaire sera formalisée en Première.
\begin{exemple}Un échantillon de \qty{2,994}{g} d'eau contient $N=2,994\times10^{-3}/(2,994\times10^{-26})=10^{23}$ molécules ; $n=10^{23}/(6,02\times10^{23})\simeq\qty{0,166}{mol}$.\end{exemple}
\essSection{4}{Expérimenter et modéliser}
\begin{experience}[Peser pour compter]Comparer un lot de rondelles identiques et un échantillon de matière : la masse moyenne d'une rondelle permet de compter indirectement le lot. Identifier l'hypothèse d'identité des objets. Transposer ensuite la démarche à des atomes ou molécules en utilisant les masses fournies, sans prétendre peser une molécule sur une balance scolaire.\end{experience}
\exercices
\rubrique{Prendre confiance}
\begin{blocLivre}{ex-6-1}\begin{exercice}[1]{Lire les indices}
Une molécule de glucose a pour formule \ce{C6H12O6}.
\begin{questions}
\item Combien d’atomes contient-elle au total ?
\item Combien d’atomes d’hydrogène contiennent 100 molécules ?
\end{questions}
\end{exercice}\end{blocLivre}
\begin{blocLivre}{ex-6-2}\begin{exercice}[1]{Masse d’une molécule}
$m_C=\qty{1,99e-26}{kg}$ et $m_O=\qty{2,66e-26}{kg}$.
\begin{questions}
\item Calculer la masse d’une molécule de \ce{CO2}.
\item Exprimer cette masse en grammes.
\end{questions}
\end{exercice}\end{blocLivre}
\begin{blocLivre}{ex-6-3}\begin{exercice}[1]{Un nombre d’atomes}
Un échantillon pur de cuivre pèse 1,06 g. Un atome de cuivre pèse \qty{1,06e-25}{kg}.
\begin{questions}
\item Convertir la masse de l’échantillon en kilogrammes.
\item Calculer le nombre d’atomes.
\end{questions}
\end{exercice}\end{blocLivre}
\begin{blocLivre}{ex-6-4}\begin{exercice}[1]{Passer à la mole}
Un échantillon contient $3,01\times10^{23}$ molécules. On adopte $N_A=\qty{6,02e23}{mol^{-1}}$.
\schema{mole}
\begin{questions}
\item Calculer la quantité de matière.
\item Combien d’entités y aurait-il dans 2,00 mol ?
\end{questions}
\end{exercice}\end{blocLivre}
\rubrique{S’entraîner}
\begin{blocLivre}{ex-6-5}\begin{exercice}[2]{Le dioxygène}
Un atome d’oxygène pèse \qty{2,66e-26}{kg}. Un échantillon de \ce{O2} pèse 0,532 g.
\begin{questions}
\item Calculer la masse d’une molécule de \ce{O2}.
\item Déterminer $N$ puis $n$ avec $N_A=\qty{6,02e23}{mol^{-1}}$.
\end{questions}
\end{exercice}\end{blocLivre}
\begin{blocLivre}{ex-6-6}\begin{exercice}[2]{Atomes et molécules}
Une quantité de 0,200 mol de \ce{CO2} est considérée.
\begin{questions}
\item Donner les quantités d’atomes de carbone et d’oxygène correspondantes.
\item Peut-on dire qu’il y a 0,200 mol d’atomes au total ?
\end{questions}
\end{exercice}\end{blocLivre}
\begin{blocLivre}{ex-6-7}\begin{exercice}[2]{Revenir à la masse}
On souhaite $2,00\times10^{22}$ molécules de masse individuelle \qty{3,00e-26}{kg}.
\begin{questions}
\item Calculer la masse de l’échantillon.
\item Exprimer cette masse dans une unité pratique pour une balance.
\end{questions}
\end{exercice}\end{blocLivre}
\begin{blocLivre}{ex-6-8}\begin{exercice}[2]{Une espèce dans un mélange}
Un échantillon de 5,00 g contient 20,0 \% en masse d’une espèce dont chaque entité pèse \qty{5,00e-26}{kg}.
\begin{questions}
\item Calculer la masse de cette espèce.
\item Calculer son nombre d’entités.
\end{questions}
\end{exercice}\end{blocLivre}
\rubrique{Raisonner et approfondir}
\begin{blocLivre}{ex-6-9}\begin{exercice}[3]{Même nombre, masses différentes}
Une molécule A pèse $3,00\times10^{-26}$ kg et une molécule B $6,00\times10^{-26}$ kg.
\begin{questions}
\item Comparer les masses de deux échantillons contenant chacun $10^{23}$ molécules.
\item Comparer leurs quantités de matière.
\end{questions}
\end{exercice}\end{blocLivre}
\begin{blocLivre}{ex-6-10}\begin{exercice}[3]{Un solide ionique}
Un échantillon contient $6,02\times10^{22}$ motifs \ce{CaCl2}. Chaque motif correspond à un ion calcium et deux ions chlorure.
\begin{questions}
\item Calculer les nombres d’ions de chaque sorte.
\item Donner leurs quantités de matière.
\end{questions}
\end{exercice}\end{blocLivre}
\begin{blocLivre}{ex-6-11}\begin{exercice}[3]{Une très petite quantité}
Une goutte de masse 0,0300 g ne contient que des molécules de masse \qty{3,00e-26}{kg}.
\begin{questions}
\item Calculer le nombre de molécules.
\item Comparer ce nombre à un milliard.
\end{questions}
\end{exercice}\end{blocLivre}
\begin{blocLivre}{ex-6-12}\begin{exercice}[3]{Détecter un facteur mille}
Un élève divise 1,80 g par $3,00\times10^{-26}$ kg et annonce $6,00\times10^{25}$ molécules.
\begin{questions}
\item Identifier son erreur.
\item Corriger $N$ puis calculer $n$.
\end{questions}
\end{exercice}\end{blocLivre}

